% detailed_project_guide.tex
% A long-form, pedagogical tutorial guide to our Phase-2 project.
% No page limit. Compile: pdflatex detailed_project_guide.tex  (run twice for the ToC/refs)
% Zero exotic dependencies -- compiles on any standard TeX Live install.
\documentclass[11pt]{article}

\usepackage[margin=1in]{geometry}
\usepackage{amsmath,amssymb}
\usepackage{booktabs}
\usepackage{xcolor}
\usepackage{enumitem}
\usepackage[hidelinks]{hyperref}

\setlist{itemsep=2pt,topsep=3pt}

% A dependency-free "callout" for intuitions and takeaways: a light rule-bounded
% quotation with a bold lead-in. Compiles anywhere; no tcolorbox/mdframed needed.
\newenvironment{takeaway}
  {\par\vspace{4pt}\noindent\begin{center}\begin{minipage}{0.94\linewidth}%
   \color{black!88}\rule{\linewidth}{0.7pt}\par\vspace{3pt}\small\itshape}
  {\par\vspace{3pt}\rule{\linewidth}{0.7pt}\end{minipage}\end{center}\vspace{4pt}}

\newcommand{\intuition}[1]{\vspace{2pt}\noindent\textbf{Intuition.}\ #1\par\vspace{2pt}}
\newcommand{\Ob}{O_0}

\title{\vspace{-1.5em}\textbf{From Hidden Confounders to Honest Answers}\\[0.3em]
\large A Ground-Up Guide to Our Confounded-POMDP Project}
\author{Deep RL Project (Phase 2) --- Sharif University of Technology\\
S.~Sayah Varg, M.~Shirinbayan, M.~Ostadmohammadi}
\date{}

\begin{document}
\maketitle

\begin{abstract}
\noindent This is a self-contained tutorial. If you know what \emph{states},
\emph{actions}, \emph{rewards}, and \emph{MDPs} are, you know enough to read every
word of it. Everything else is built up from scratch --- analogies first, equations
second. We explain the problem our project attacks (learning from decision logs
poisoned by a hidden variable), the elegant 2024 theory we implemented to fix it,
the place where that theory quietly explodes when you actually run it, the repair we
invented, and --- being honest --- the sobering limits of both the repair and the
whole approach. Read it top to bottom like a story; each section sets up the next.
\end{abstract}

\tableofcontents
\vspace{1em}

% =====================================================================
\section*{The 60-Second Version}
\addcontentsline{toc}{section}{The 60-Second Version}

We do \textbf{offline reinforcement learning}: we have a big log of decisions someone
already made (say, doctors treating patients) and we want to judge \emph{new}
strategies without ever running them for real. The catch is that the log is
\textbf{poisoned by a hidden variable} --- something the historical decision-maker
could see and react to, but which was never written down. That hidden variable makes
the naive way of scoring strategies \emph{systematically wrong}, and no amount of
extra data fixes it. A 2024 paper (Hong, Qi \& Xu) proposed an elegant cure using two
objects called \textbf{bridge functions} plus a special ``before-the-decision''
measurement. It was pure theory. \textbf{We built it, ran it, and stress-tested it},
found where the math explodes in practice, invented a fix, and honestly mapped where
the fix (and the whole method) helps versus where it loses to a dumb baseline.

% =====================================================================
\section{The Intuitive Problem: Confounding and Partial Observability}

\subsection{The comfortable world: a normal MDP}
In a standard \textbf{Markov Decision Process (MDP)} you see the full \textbf{state}
$s$, take an \textbf{action} $a$, get a \textbf{reward} $r$, and the world moves to
$s'$. ``Markov'' means the state holds everything relevant: given $s$ and $a$, the
future ignores the past. Here, offline RL is easy in principle: collect a log of
$(s,a,r,s')$, average ``after state $s$ and action $a$, what happened,'' and you can
score any strategy. The quiet assumption baking this cake is that \textbf{you can see
the whole state}. That is exactly what breaks in reality.

\subsection{Two things go wrong at once}
\textbf{Partial observability.} You rarely see the true state $s$. In an ICU you see
\emph{vital signs} $O$ (the observation), but the true physiological state $s$ is
hidden --- vitals are a blurry shadow of it. A system where you see only observations
of a hidden state is a \textbf{POMDP} (Partially Observable MDP).

\textbf{Confounding.} Here is the killer. The doctor who generated your log
\emph{could see more than you can}: training, intuition, an unrecorded sense of how
sick the patient ``really'' was. They used that hidden information to \emph{choose the
treatment}, and the same hidden information also drove the \emph{outcome}. A hidden
variable that touches both the action and the outcome is a \textbf{confounder}. It is
the single most dangerous thing in causal inference.

\subsection{An analogy that makes it click}
You study a company's sales logs to decide whether sending a salesperson to a client
is worth it. You notice: \emph{visited clients almost always buy}. Naive conclusion:
\emph{visiting works, send salespeople to everyone!} But a hidden variable the sales
manager saw and you cannot: \textbf{which clients were already eager to buy}. The
manager sent salespeople \emph{precisely} to the clients about to buy anyway. The
visits did not cause the sales --- the hidden ``eagerness'' caused both the visit and
the sale. Roll out ``visit everyone'' and you will be disappointed, because you now
visit un-eager clients too. That hidden eagerness is the confounder; it made your log
say visiting is magic when it is not. This is exactly our situation, with doctors,
treatments, and hidden patient severity.

\subsection{Why naive evaluation is \emph{systematically} biased (not just noisy)}
This is the point people underestimate. When you score a new strategy naively ---
treating the observation $O$ as the state and averaging log outcomes --- you do not
get a \emph{noisy} answer that averages out with more data. You get a \textbf{biased}
answer that stays wrong \emph{no matter how much data you collect}. The reason: the
log's actions are correlated with the hidden confounder, and so are the outcomes, so
your average silently mixes ``the effect of the action'' with ``the effect of the
hidden thing that made this action get chosen.'' More data just estimates \emph{the
wrong number} more precisely.

In our experiments this is concrete. On our controllable test environment (the
``toy''), the naive evaluator over-scores policies by $+0.14$ to $+0.33$ in value
units ($8\%$ to $22\%$ relative), and this bias is \textbf{flat as $N$ grows}: at
$N=5{,}000$ and $N=50{,}000$ it barely moves. That flatness is the fingerprint of
\emph{bias}, not variance --- variance shrinks with data, bias does not.

\begin{takeaway}
Offline RL breaks when a hidden variable drove both the logged actions and the
outcomes. The naive ``pretend the observation is the state'' evaluator is then
\emph{systematically} biased, and --- unlike noise --- that bias never washes out
with more data. That is the enemy.
\end{takeaway}

% =====================================================================
\section{The Anchor's Core Solution: The Dual-Bridge Idea}

Our project builds on \textbf{Hong, Qi \& Xu, ``Model-based RL for Confounded
POMDPs'' (ICML 2024)} --- the ``anchor paper.'' It gives a beautiful way to
un-poison the log. We build it up in plain language before any equation.

\subsection{The secret ingredient: a ``negative control'' taken \emph{before} the decision}
The trick starts with an extra measurement we call $\Ob$: a reading taken
\emph{before the first action is chosen} --- a snapshot at the very start.

\intuition{Because $\Ob$ happens \emph{before} the doctor acts, it is a clean
``echo'' of the hidden state that is \textbf{not contaminated by the decision}. It is
like a \textbf{fingerprint of the hidden state, taken at the door before the crime}.
It does not reveal the hidden state directly (it is noisy), but it is correlated with
it \emph{and} untainted --- exactly what lets us correct for a variable we never
actually observe.}

In causal-inference language $\Ob$ is a \textbf{negative control} (or \emph{proxy}).
Its formal job (the paper's Assumption~3.1) is: \emph{given the true hidden state and
the action, $\Ob$ is independent of the later observations, next-states, and
rewards}. In words: once you know the hidden truth, the early fingerprint tells you
nothing new. All of $\Ob$'s value is that it is an untainted window onto the hidden
state --- nothing sneaky.

\subsection{What a ``bridge function'' actually is}
Now the centerpiece: two \textbf{bridge functions}, $b_R$ (reward bridge) and $b_D$
(dynamics bridge).

\begin{takeaway}
A bridge function is a \textbf{re-expression of a hidden-state quantity purely in
terms of things you can see}, built so that its average --- \emph{viewed through the
lens of the negative control $\Ob$} --- matches a quantity you can measure from data.
\end{takeaway}

\intuition{You want to know something about the hidden state (e.g.\ the true reward a
policy would earn), but you only see observations. A bridge function is a
\textbf{corrective lens}: a function of the \emph{observed} variables shaped so that
looking at it \emph{through the $\Ob$ window} reproduces the real conditional
probabilities in the data. The deep theorem says: if a function of observables
reproduces the observable data when viewed through $\Ob$, then that same function
reproduces the \emph{hidden-state} quantity you actually wanted. The $\Ob$ window ties
the two together.}

There are two bridges because there are two hidden-state quantities to rebuild:
\begin{itemize}
  \item \textbf{$b_R$ --- the reward bridge:} recovers the reward model (what rewards
  a policy would really earn).
  \item \textbf{$b_D$ --- the dynamics bridge:} recovers the transition model (how the
  hidden state would really evolve).
\end{itemize}
Together they rebuild the \emph{entire} confounded-POMDP machinery from the poisoned
log (this is why the method is ``model-based'' and ``dual-bridge''). Chain them and
you can evaluate \emph{any} policy cleanly, with confounding subtracted out.

\subsection{From infinite theory to finite matrices: where approximation sneaks in}
The paper is written for \emph{infinite populations and integral equations}. Reality
gives a \emph{finite pile of trajectories}. The gap is where every practical
difficulty lives. The bridge's defining ``matching condition'' is really a
\textbf{Fredholm integral equation of the first kind} --- ``solve for the function
$b$ inside an averaging operator'':
\begin{equation}
\underbrace{\mathbb{E}\big[\,b(W_t, y)\;\big|\;X_t\,\big]}_{\text{average of }b\text{ you can see, given the }\Ob\text{-window}}
\;=\;\underbrace{p(y\mid X_t)}_{\text{a probability you can measure}},
\end{equation}
for every index $y$, where $W_t=(A_t,O_t)$ are the observed variables and
$X_t=(A_t,H_{t-1},\Ob)$ is the conditioning that includes the negative control.

\intuition{Solving for $b$ means \textbf{inverting an averaging operator}. But
averaging \emph{smooths things and throws information away}, so un-averaging it is
\textbf{ill-posed}: tiny noise wiggles in the measured right-hand side get
\emph{amplified} into big wiggles in the recovered $b$. This ill-posedness is the
villain of Section~3.}

The paper's estimator, which we implemented, handles this in \textbf{two stages}:
\begin{itemize}
  \item \textbf{Stage~1 (learn the lens).} From part of the data, estimate the
  averaging operator itself (a \emph{conditional mean embedding}). In our tabular
  setting this is literally \textbf{counting}: build tables of conditional
  frequencies. Call the result the ``lens matrix'' $\Gamma$.
  \item \textbf{Stage~2 (solve for the bridge).} From the \emph{other} part of the
  data, solve the matching equation using the Stage-1 lens. In closed form:
  \begin{equation}
  \alpha \;=\; (G + N_2\,\lambda_2\,I)^{-1}\,\hat p,
  \qquad G=(\Gamma^\top K_W\,\Gamma)\odot K_{Y''},
  \label{eq:stage2}
  \end{equation}
  where $G$ is built from the lens and the data, $\hat p$ is the measured right-hand
  side, $\odot$ is an element-wise product, and $\alpha$ are the coefficients that
  \emph{are} the estimated bridge.
\end{itemize}
The only genuinely new piece in \eqref{eq:stage2} is the term
$+\,N_2\lambda_2 I$: the \textbf{regularizer}.

\intuition{Because inverting an ill-posed operator amplifies noise, we deliberately
add a little ``stiffness'' ($\lambda_1$ in Stage~1, $\lambda_2$ in Stage~2) that
stops the solution from chasing noise --- at the cost of a little bias. It is the
universal \textbf{bias--variance knob}: turn $\lambda$ up for a stable-but-slightly-wrong
answer; turn it down for a less-biased but wildly noisy one.}

\subsection{What our hyperparameter calibration actually did}
The paper prescribes \emph{theoretical schedules} like $\lambda\propto N^{-1/2}$,
derived for the \emph{infinite-dimensional kernel} version. We implemented the
\emph{tabular} version --- and hit our first real finding.

\begin{takeaway}
The paper's theoretical $\lambda$ schedules do \textbf{not} transfer to the tabular
setting. Plugging in $\lambda_1\propto N^{-1/2}$ injected $\sqrt{N_1}$ fake
``pseudo-counts'' into our count tables, over-shrinking the Stage-1 lens by roughly
$80\%$ by the third time-step. The theory's stiffness was calibrated for a different
object and was far too stiff for ours.
\end{takeaway}

So we did the honest thing: a \textbf{grid calibration} --- sweep the knobs over a
small grid, across three seeds, and \emph{measure} which settings minimize error. The
winners:
\begin{itemize}
  \item \textbf{Data split $N_1:N_2 = 0.7:0.3$} --- give Stage~1 (learning the lens
  over a growing history) more data. This confirmed the paper's qualitative ``Stage~1
  needs more'' intuition.
  \item $\boldsymbol{\lambda_1 = 1/N_1}$ --- a gentle, parametric-rate stiffness.
  \item $\boldsymbol{\lambda_2 = 0.03/\sqrt{N_2}}$ --- crucially, this decays
  \emph{slower} than $1/N_2$. It must, because $\lambda_2$'s job is to
  \textbf{dominate the noise leaking out of Stage~1} into the weakly-identified
  directions of the bridge. Too fast, and that leaked noise wins.
\end{itemize}
A subtler lesson that matters in Section~3: the obvious ``is my matrix healthy?''
check is its \textbf{condition number} --- and here it is a \textbf{liar}. Because of
the null-space structure (Section~3), it stays flat and reassuring even as the
estimate quietly collapses. The honest health-check is the smallest \emph{signal}
eigenvalue $\hat\sigma_2$ --- the weakest direction that still carries real
information. We verified this by degrading $\Ob$: the true Stage-1 error blew up
$10\times$ while the condition number \emph{did not budge}. So all our stability
monitoring watches $\hat\sigma_2$, not the condition number.

\begin{takeaway}
A ``bridge'' is an observable function whose $\Ob$-windowed average reproduces the
data --- and, by a deep theorem, thereby reproduces the hidden-state model you
wanted. Estimating it from finite data is a two-stage, ill-posed matrix inversion,
tamed by regularizers $\lambda_1,\lambda_2$ whose right values we had to \emph{measure},
not inherit from the theory.
\end{takeaway}

% =====================================================================
\section{The Vanilla Breakdown and Our Innovation}

This is the most interesting part: where the beautiful theory, actually run,
\textbf{detonates} --- and what we did about it.

\subsection{The plan: ``pessimism''}
Once you can estimate the model, you want to \emph{pick a good policy}. But your
bridges are noisy, so if you just trust them and pick the best-looking policy, you
fall for one that got \emph{lucky} in your noisy estimate. The standard defense is
\textbf{pessimism}: score a policy by a cautious \emph{lower bound}, and pick the
policy whose \emph{cautious} score is highest. This ``be paranoid about what you can't
pin down'' principle is provably right for offline RL.

The anchor paper draws, around the estimate $\hat b$, a \textbf{confidence region} ---
a fuzzy ball of ``bridges that fit the data almost as well as $\hat b$.'' Because the
underlying error is exactly quadratic, this region is a clean \textbf{ellipsoid}. The
pessimistic value is then
\[
V_{\text{low}}(\pi) \;=\; \min_{\,b\ \in\ \text{confidence ellipsoid}}\; V(\pi, b),
\]
``the smallest value of policy $\pi$ you can get as the bridges range over the whole
plausible region.'' Minimizing a linear value over an ellipsoid has a gorgeous
closed form --- the \textbf{support function}:
\begin{equation}
V_{\text{low}} \;=\; V(\hat b)\;-\;\sqrt{\xi}\;\lVert g\rVert_{H^{-1}},
\label{eq:supp}
\end{equation}
where $\xi$ is the ellipsoid's size, $g$ is the value's gradient, and
$\lVert\cdot\rVert_{H^{-1}}$ measures length in the ellipsoid's own stretched
geometry. In words: \emph{start at the estimated value, then subtract a penalty for
how much the value could wiggle within the plausible region}. Beautiful. So what
goes wrong?

\subsection{The detonation: why vanilla pessimism yields $-1755$}
Here is the crux, and it is a genuinely subtle geometric fact.

Recall from Section~1 that \textbf{the observation is richer than the hidden state}
($|O| > |S|$). In our benchmark, $720$ observable ``core'' states stand in for a
$2$-valued hidden variable; even in the tiny toy there are $3$ observations for $2$
hidden states. That mismatch means the data-fitting matrix has a \textbf{null space}
--- directions in ``bridge space'' the data says \emph{nothing} about. You can move
the bridge along a null direction and the data-fit \emph{does not change at all},
because those directions describe patterns in the rich observation with no
counterpart in the poor hidden state.

Now look at the ellipsoid. Its ``stiffness'' in a direction is how much moving there
hurts your fit. In the null directions, moving hurts \emph{only} through the tiny
regularizer $\lambda_2$. So the ellipsoid is \textbf{catastrophically elongated} along
the null space --- not a round ball, but a \textbf{needle} stretching almost
infinitely into the directions where the data is blind. Concretely, the pessimism may
travel a distance $\lVert d\rVert=\sqrt{\xi/\lambda_2}$ along a null direction, and
since $\lambda_2$ is tiny, that distance is \emph{enormous}.

The vanilla recipe minimizes the value \textbf{jointly over all bridge blocks at
once}. Let a joint minimization loose on a needle stretching to the horizon in the
blind directions, and it \emph{runs to the horizon}: it finds these free,
data-invisible directions, pushes the bridges way out along them, and the multi-step
value chain amplifies the excursion into nonsense. We measured exactly this:

\begin{center}\small
\begin{tabular}{@{}lccccc@{}}
\toprule
Ellipsoid size $c$ & $0.1$ & $0.3$ & $1.0$ & $3.0$ & $10.0$\\
\midrule
$V_{\text{low}}$ (best policy) & $-1.12$ & $-10.4$ & $-61.6$ & $-301.7$ & $\mathbf{-1755}$\\
\bottomrule
\end{tabular}
\end{center}

True policy values live between about $+1.5$ and $+2.2$, so a ``cautious lower
bound'' of $-1.12$ is already \emph{physically impossible}, and $-1755$ is
\textbf{deranged}. It is so busy being paranoid about directions the data cannot even
see that it produces numbers with no connection to reality and ``selects'' policies
essentially at random.

\begin{takeaway}
\textbf{One-sentence diagnosis:} the paper's pessimism spends all its paranoia on the
empty, data-blind \emph{null space} instead of on the directions where the data
actually has something to say.
\end{takeaway}

\subsection{The innovation: Signal-Projected Pessimism}
Once you \emph{see} the problem, the fix writes itself. If the disaster comes from the
pessimism wandering into the null space (the empty noise directions), then
\textbf{do not let it go there}: confine the whole computation to the \textbf{signal
subspace} --- the directions the data actually informs.

How do you find the signal subspace? With the workhorse of applied linear algebra:
the \textbf{SVD / eigen-decomposition}. Decompose the data's design matrix into
directions ranked by how much real information (signal) they carry. Top directions
with healthy eigenvalues are \textbf{signal} (the data speaks clearly). Bottom
directions with near-zero eigenvalues are the \textbf{null space} (the data is
silent). An \emph{eigengap} rule draws the line. We then build an orthonormal basis
$U$ for the signal directions and run the \emph{same} support-function pessimism, but
\textbf{restricted to $U$}:
\begin{equation}
b_{\min} \;=\; \hat b \;-\; \sqrt{\xi}\;\frac{U\,H_{\mathrm{sub}}^{-1}\,U^\top g}
{\lVert U^\top g\rVert_{H_{\mathrm{sub}}^{-1}}},
\qquad H_{\mathrm{sub}} = U^\top H\,U.
\end{equation}

\intuition{Project everything onto the signal subspace first, then be pessimistic
only within it. The pessimism can no longer escape into the blind directions because
we walled them off. The needle becomes a sensible ellipsoid again.}

\textbf{And on the surface it works, beautifully.} The projected pessimistic value
stays \emph{informative}: for the best policy it is $V_{\text{low}}=1.35$ (sane,
in-range) instead of $-1.12$. And it \textbf{picks the right policy}:
\textbf{$0.000$ selection regret}, across every random seed. Where vanilla chose
garbage, the projected version reliably identifies the truly best policy.

\subsection{The honest twist: a heuristic, not a guaranteed cure}
Here is the part we refuse to sugar-coat, because being honest here is the most
valuable thing the project did.

Walling off the null space \emph{assumes the true bridge lives inside the signal
subspace}. We \emph{checked} that assumption by measuring how much of the true bridge
leaks \emph{out} of the signal subspace. The answer: about \textbf{21\%} of its length
(for one bridge, the leaked-out part had length $0.38$ while the whole bridge had
length $1.80$).

That leakage is fatal to the \emph{guarantee}. If the truth is $21\%$ outside the
region you search, your ``confidence region'' \textbf{does not contain the truth} ---
its honest coverage is \emph{zero}. A pessimistic bound is only trustworthy if the
truth is in the set you minimize over; ours is not. So $V_{\text{low}}\le
V_{\text{true}}$ held in all our runs \emph{by luck of sign-alignment}, not because
the geometry guarantees it.

\begin{takeaway}
There is a fundamental \textbf{tension between coverage and informativeness}. The
vanilla method \emph{covers} the truth but is \emph{uninformative} (blows up to
$-1755$, picks garbage). Our projected method is \emph{informative} (sane numbers,
$0.000$ regret) but does \emph{not} cover the truth (walls off $21\%$ of it). You
cannot have both with this estimator. Signal-Projected Pessimism is an excellent,
practically-useful \textbf{heuristic} --- but it is \emph{not} a mathematically
guaranteed pessimistic bound, and we say so plainly.
\end{takeaway}

The real intellectual product here is not ``we fixed the paper,'' but ``we found
precisely \emph{why} the paper's beautiful pessimism is unusable as written, we made
it usable, and we mapped the exact price of doing so.''

% =====================================================================
\section{The Empirical Verdict}

Theory and analogies are nice, but this was fundamentally an
\emph{engineering-and-measurement} project. Here is what we actually observed, in
order of increasing humility.

\subsection{The math is implemented correctly (the $10^{-15}$ handshake)}
Before trusting any result, prove the code computes what the equations say. We used
\textbf{oracle checks} --- comparisons against a ground truth computed exactly in a
small environment:
\begin{itemize}
  \item The value chain matches exact dynamic programming to $\mathbf{4.4\times
  10^{-16}}$. Feeding the \emph{true} bridges into our evaluator reproduces the
  exactly-correct value to machine precision (essentially the smallest gap a 64-bit
  computer can represent). Our tensor contractions \emph{are} the equations, bit for
  bit.
  \item The estimator hits the right target to $\mathbf{3.0\times10^{-15}}$: in the
  infinite-data limit our bridge solver lands on the exact ideal solution.
  \item Two independent derivations (a ``dual'' and a ``primal'' form of
  Eq.~\eqref{eq:stage2}) agree to $\mathbf{2.4\times10^{-15}}$ --- a strong internal
  consistency proof.
\end{itemize}
These $10^{-15}$-level agreements are the bedrock: every \emph{interesting} result
later is a fact about the \emph{method}, not a bug in the \emph{code}.

\intuition{A humbling aside that captures the whole project's spirit: while building a
comparison baseline, our first version was a perfectly reasonable-looking derivation
that the oracle immediately exposed as \emph{systematically wrong} (it gave $1.61$ for
a policy whose true value is $1.865$, and the error did not shrink with data). Only
running it against the exact answer caught it. The lesson we kept relearning:
\textbf{reasoning proposes, the oracle disposes}.}

\subsection{More data $\rightarrow$ smoothly less error}
On the toy, where we can measure bridge-recovery error against the true bridges, the
error drops \textbf{monotonically at every step} as we add data:

\begin{center}
\begin{tabular}{@{}lcccc@{}}
\toprule
Trajectories $N$ & $1{,}000$ & $5{,}000$ & $10{,}000$ & $20{,}000$\\
\midrule
Reward-bridge error $\mathrm{err}_{b_R}$ & $1.360$ & $0.775$ & $0.581$ & $\mathbf{0.459}$\\
Dynamics-bridge error $\mathrm{err}_{b_D}$ & $0.732$ & $0.422$ & $0.312$ & $\mathbf{0.223}$\\
\bottomrule
\end{tabular}
\end{center}

Fitting a line on a log--log plot gives decay rates of about $N^{-0.36}$ and
$N^{-0.39}$ --- steady, honest statistical convergence. The payoff: our de-biased
evaluator beats the naive (confounding-blind) baseline in \textbf{12 out of 12} test
cases on the toy. \emph{Where confounding is strong and identification is clean, the
method does exactly what it promised: it removes the bias that more data alone never
could.} This is the success story.

\subsection{The boundary condition: where it loses to a dumb baseline}
Now the humbling and most scientifically important result. Scaling to the full
benchmark ($720$ observable states, a realistic size), the story \textbf{reversed}.
We ran a comprehensive sweep --- five confounding strengths $\kappa$ from $0$ (none)
to $1$ (maximal), three data sizes up to $60{,}000$, three seeds, four methods --- and
measured every method against the true policy values.

\begin{takeaway}
The naive, confounding-blind baseline \textbf{beats} our sophisticated de-biasing
method almost everywhere at benchmark scale. The plug-in's error ($0.006$--$0.12$) is
comparable to or \emph{larger} than the structural bias it was built to remove
($0.004$--$0.011$).
\end{takeaway}

\intuition{Recall the two enemies: \emph{bias} (from confounding) and \emph{variance}
(from finite data). The bridge method kills the \emph{bias} --- but killing it
requires inverting ill-posed matrices, and \textbf{inverting ill-posed matrices
amplifies noise}, i.e.\ it \emph{injects variance}. At benchmark scale each little
$2\times2$ correction is estimated from a handful of samples, so the variance it
injects is \emph{bigger than the bias it removes}. You pay more in noise than you gain
in bias reduction. The cure is worse than the disease.}

The most telling data point: the method's \emph{worst} performance is at
$\boldsymbol{\kappa=0}$, where there is \emph{no confounding at all} --- up to
$\mathbf{23}$--$\mathbf{46\times}$ worse than naive. A correction tool that damages
the answer when there is nothing to correct is telling you loudly that its problem is
\emph{variance}, not bias.

We tried to fix this honestly. A \textbf{James--Stein shrinkage} upgrade (pool each
noisy little matrix toward a stable, confounding-aware group average) genuinely helped
where it should: under strong confounding it \textbf{halved} the error
($0.22\rightarrow0.08$) and cut run-to-run variance $\mathbf{6\times}$. A real,
defensible improvement --- but \emph{not enough to overturn the boundary}; shrinkage
narrows the gap to naive without closing it.

We also built a \textbf{model-free} competitor (estimate value functions directly
instead of the model). It was worse and \emph{flatter} still --- its error
\emph{did not improve with more data at all}. The reason is instructive: our negative
control $\Ob$ is only a \emph{binary} fingerprint, too little information to pin down
a value function over $720$ observations. It is \textbf{under-identified} --- like
reconstructing a detailed photograph from a single bit. (Reassuringly, on the toy
where identification is clean, this same method \emph{is} consistent --- so its
benchmark failure is a genuine identification limit, not a bug.)

\subsection{The honest bottom line}
\begin{itemize}
  \item \textbf{Model-based beats model-free} at benchmark scale (the model-free
  method is starved by the tiny proxy).
  \item \textbf{But both lose to naive} at realistic sample sizes, because estimation
  \emph{variance} overwhelms the confounding \emph{bias} they are designed to remove.
  \item \textbf{Shrinkage helps under strong confounding} but does not change the
  verdict.
  \item \textbf{The method shines exactly where the theory promised} --- strong
  confounding, clean identification, enough data (our toy) --- and struggles exactly
  where practice bites (large observation spaces, weak proxies, limited data).
\end{itemize}

\begin{takeaway}
The through-line of the entire project, in one sentence: \textbf{at realistic scale,
variance --- not identification --- is the wall.} The anchor paper solved the
identification problem elegantly; what it (understandably, being pure theory) never
had to confront is that solving identification by matrix inversion re-introduces a
variance problem that can be worse than the bias you started with.
\end{takeaway}

% =====================================================================
\section*{Epilogue: What This Project Taught Us}
\addcontentsline{toc}{section}{Epilogue}

\begin{enumerate}
  \item \textbf{A gorgeous theory can be silently unusable, and only running it
  reveals where.} The $-1755$ blow-up, the null-space geometry, the $21\%$ leak, the
  benchmark reversal --- none are visible in the paper. They appeared only when exact
  code met exact oracles. Every real finding came from \emph{measurement}, not from
  thinking harder.
  \item \textbf{Honesty about a heuristic beats a fake theorem.} We could have called
  Signal-Projected Pessimism a ``repair'' and moved on. Instead we measured the $21\%$
  leak and reported the coverage/informativeness tension. That negative result is more
  useful --- and more true --- than a victory lap.
  \item \textbf{The bias--variance trade-off is the deepest force in the room.}
  Confounding is a bias problem; the elegant fix is a variance-injecting inversion; at
  scale the variance wins. Understanding \emph{that}, and mapping exactly where the
  crossover happens, is the most valuable thing to hand to whoever builds on this next.
\end{enumerate}

\noindent If you read this far, you now understand --- from first principles --- a
genuine research-grade problem in causal reinforcement learning, an elegant proposed
solution, exactly how and why it breaks in practice, a clever-but-honestly-limited
fix, and the sober empirical reality of when it helps and when it does not. That full
arc --- the promise, the breakdown, the repair, and the honest verdict --- is the
real education.

\end{document}
